\begin{statementbox}
\textbf{条件与对象。} $p>0$，直线 $l:Ax+By+C=0$，$A^2+B^2>0$。
交点记为 $P(x_1,y_1),Q(x_2,y_2)$；相切时按重合端点计。
以下韦达关系解释为实交点坐标时，\textbf{必须先确认实交点存在}。

\textbf{一、保留横坐标的四种固定形式。}

\textbf{向右：$y^2=2px$，二次韦达需 $A\ne0$。}
\[
 A^2x^2+(2AC-2B^2p)x+C^2=0,
\]
\[
 x_1+x_2=\frac{2B^2p-2AC}{A^2},\qquad x_1x_2=\frac{C^2}{A^2}.
\]
\textbf{向左：$y^2=-2px$，二次韦达需 $A\ne0$。}
\[
 A^2x^2+(2AC+2B^2p)x+C^2=0,
\]
\[
 x_1+x_2=-\frac{2B^2p+2AC}{A^2},\qquad x_1x_2=\frac{C^2}{A^2}.
\]
\textbf{向上：$x^2=2py$，二次韦达需 $B\ne0$。}
\[
 Bx^2+2Apx+2Cp=0,
 \qquad x_1+x_2=-\frac{2Ap}{B},\quad x_1x_2=\frac{2Cp}{B}.
\]
\textbf{向下：$x^2=-2py$，二次韦达需 $B\ne0$。}
\[
 Bx^2-2Apx-2Cp=0,
 \qquad x_1+x_2=\frac{2Ap}{B},\quad x_1x_2=-\frac{2Cp}{B}.
\]
\textbf{左右开口有一个额外提醒。}
若 $B\ne0$，横坐标根可唯一恢复 $y=-(Ax+C)/B$。
若 $B=0$，上式只剩 $(Ax+C)^2=0$，\textbf{单看这个重根不能确认实交点存在，也不能判断相切}。
此时改用下一页的纵坐标方程，覆盖竖直弦、顶点切线和无交点三种情况。
\end{statementbox}
\clearpage
\begin{statementbox}
\textbf{二、先选能唯一恢复交点的坐标。}
令 $s=1$ 或 $-1$，分别表示沿正轴或负轴开口；$p$ 始终为正。

\textbf{左右开口 $y^2=2spx$：优先消去 $x$。} 当 $A\ne0$，有
\[
 Ay^2+2spBy+2spC=0,\qquad
 y_1+y_2=-\frac{2spB}{A},\quad y_1y_2=\frac{2spC}{A}.
\]
记 $H=p^2B^2-2spAC$，则 $\Delta_y=4H$。
当 $A=0$ 时 $B\ne0$，直线平行于抛物线轴，\textbf{恰有一个交点且非相切}：
\[
 \left(\frac{sC^2}{2pB^2},-\frac CB\right).
\]
\textbf{上下开口 $x^2=2spy$：优先消去 $y$。} 当 $B\ne0$，有
\[
 Bx^2+2spAx+2spC=0,\qquad
 x_1+x_2=-\frac{2spA}{B},\quad x_1x_2=\frac{2spC}{B}.
\]
记 $K=p^2A^2-2spBC$，则 $\Delta_x=4K$。
当 $B=0$ 时 $A\ne0$，直线平行于抛物线轴，\textbf{恰有一个交点且非相切}：
\[
 \left(-\frac CA,\frac{sC^2}{2pA^2}\right).
\]
上述\textbf{二次项非零时}，$H$ 或 $K$ 的正、零、负，分别对应两个不同实交点、相切、无实交点。

\textbf{三、中点与弦长。} 当有两个不同实交点时，记中点为 $S$：
\[
\begin{array}{c|c|c}
 &S&PQ\\[3pt]\hline
 y^2=2spx&\left(\dfrac{spB^2-AC}{A^2},-\dfrac{spB}{A}\right)
 &\dfrac{2\sqrt{(A^2+B^2)H}}{A^2}\\[10pt]
 x^2=2spy&\left(-\dfrac{spA}{B},\dfrac{spA^2-BC}{B^2}\right)
 &\dfrac{2\sqrt{(A^2+B^2)K}}{B^2}
\end{array}
\]
第一行需 $A\ne0,H>0$；第二行需 $B\ne0,K>0$。
这些公式包含竖直弦、水平弦。相切时中点公式给出切点，长度为 $0$，不再是一条非零弦。
\end{statementbox}
\clearpage
